% ============================================================================
% REFERENCES
% ============================================================================

\chapter*{منابع و مآخذ}
\addcontentsline{toc}{chapter}{منابع و مآخذ}

\begin{LTR}
\begin{thebibliography}{99}

\bibitem{ref1} Yin, H., Xu, Y., Zhang, Y., \& Wu, J. (2024). Image compression encryption algorithm combining two-dimensional modular hyperchaotic map and compressed sensing. \textit{Physica Scripta}, 99(10), 105288.

\bibitem{ref2} Wang, X. Y., Zhang, Y. Q., \& Zhao, Y. Y. (2015). A novel image encryption scheme based on 2-D logistic map and DNA sequence operations. \textit{Nonlinear Dynamics}, 82, 1269-1280.

\bibitem{ref3} Zhang, B., \& Liu, L. (2023). Chaos-based image encryption: Review, application, and challenges. \textit{Mathematics}, 11(11), 2585.

\bibitem{ref4} Hua, Z., \& Zhou, Y. (2019). Exponential chaotic model for generating robust chaos. \textit{IEEE transactions on systems, man, and cybernetics: systems}, 51(6), 3713-3724.

\bibitem{ref5} Li, L. (2024). A novel chaotic map application in image encryption algorithm. \textit{Expert Systems with Applications}, 252(Part B), 124316.

\bibitem{ref6} Alvarez, M Shakir, D. A. Q., Salim, A., Abd Al-Rahman, S. Q., \& Sagheer, A. M. (2023). Image encryption using lorenz chaotic system. \textit{Journal of Techniques}, 5(1), 122-128.

\bibitem{ref7} Khodadadi, H., \& Zandvakili, A. (2019). A New Method for Encryption of Color Images based on Combination of Chaotic Systems. \textit{Journal of AI and Data Mining}, 7(3), 377-383.

\bibitem{ref8} Wang, X., \& Li, Y. (2021). Chaotic image encryption algorithm based on hybrid multi-objective particle swarm optimization and DNA sequence. \textit{Optics and Lasers in Engineering}, 137, 106393.

\bibitem{ref9} Mamat, A. R., Mohamed, M. A., Abidin, A. F. A., Mohamed, R. R., Sambas, A., Rusyn, V., ... \& Markovych, B. M. (2024). Color image encryption using chaotic-based cryptosystem. \textit{Mathematical Modeling and Computing}, 11(3).

\bibitem{ref10} Li, Y., Wang, X., \& Jin, J. (2020). A new one-dimensional chaotic system with applications in image encryption. \textit{Chaos, Solitons \& Fractals}, 139, 110102.

\bibitem{ref12} Ge, B., Chen, X., Chen, G., \& Shen, Z. (2021). Secure and fast image encryption algorithm using hyper-chaos-based key generator and vector operation. \textit{IEEE Access}, 9, 137635-137654.

\bibitem{ref13} Hosny, K., El-Nabawy, Y., Elshewey, A. M., Alhammad, S., et al. (2024). New method of colour image encryption using triple chaotic maps. \textit{IET Image Processing}, 18(12), 3262--3276.

\bibitem{ref14} Chen, D., Zhu, Z., \& Yang, G. (2008). An improved image encryption algorithm based on chaos. \textit{In The 9th International Conference for Young Computer Scientists}, 2792-2796.

\bibitem{ref15} Lawande, Q. V., Ivan, B. R., \& Dhodapkar, S. D. (2005). Chaos based cryptography: a new approach to secure communications. \textit{BARC newsletter}, 258, 258-2.

\bibitem{ref16} Celikovsky, S., \& Chen, G. (2005). On the generalized Lorenz canonical form. \textit{Chaos, Solitons \& Fractals}, 26(5), 1271-1276.

\bibitem{ref17} Zhou, T., Tang, Y., \& Chen, G. (2004). Chen's attractor exists. \textit{International Journal of Bifurcation and Chaos}, 14(09), 3167-3177.

\bibitem{ref18} Hsiao, H. Y. (2005). \textit{A Study of Chaotic Image Encryption}. Master of Science Thesis, Institute of Electrical Engineering, Chung Yuan Christian University, Taiwan.

\bibitem{ref19} Li, S. (2004). When chaos meets computers. \textit{arXiv preprint nlin/0405038}.

\bibitem{ref20} Strogatz, S. H. (1996). \textit{Nonlinear dynamics and chaos}. Westview Press.

\bibitem{ref21} Stojanovski, T., Pihl, J., \& Kocarev, L. (2001). Chaos-based random number generators. Part II: practical realization. \textit{IEEE Transactions on Circuits and Systems I}, 48(3), 382-385.

\bibitem{ref22} Li, S., Chen, G., Cheung, A., Bhargav, B., \& Lo, K. T. (2007). On the design of perceptual MPEG Video encryption algorithms. \textit{IEEE Trans Circuits Syst Video Technol}, 17(2), 214-223.

\bibitem{ref23} Lorenz, E. N. (1963). Deterministic nonperiodic flow. \textit{Journal of the atmospheric sciences}, 20(2), 130-141.

\bibitem{ref24} Mandelbrot, B. B. (1977). \textit{The fractal geometry of nature}. WH freeman.

\bibitem{ref25} Feigenbaum, M. J. (1978). Quantitative universality for a class of nonlinear transformations. \textit{Journal of statistical physics}, 19(1), 25-52.

\bibitem{ref27} Abuturab, M. R. (2015). An asymmetric single-channel color image encryption based on Hartley transform and gyrator transform. \textit{Opt Lasers Eng}, 69, 49-57.

\bibitem{ref28} Chen, J. X., Zhu, Z. L., Liu, Z., Fu, C., Zhang, L. B., \& Yu, H. (2014). A novel double-image encryption scheme based on cross-image pixel scrambling in gyrator domains. \textit{Opt Express}, 22(6), 7349-61.

\bibitem{ref29} Joshi, M., Shakher, C., \& Singh, K. (2008). Image encryption and decryption using fractional Fourier transform and radial Hilbert transform. \textit{Opt Lasers Eng}, 46(7), 522-6.

\bibitem{ref30} Liu, Z., Dai, J. M., Sun, X. G., \& Liu, S. T. (2010). Color image encryption by using the rotation of color vector in Hartley transform domains. \textit{Opt Lasers Eng}, 48(7), 800-5.

\end{thebibliography}
\end{LTR}

\begin{RTL}
\begin{thebibliography}{99}
\setcounter{enumiv}{30}

\bibitem{ref11} صادقیان، سعیده (۱۴۰۰). رمزنگاری تصویر مبتنی بر آشوب با ترکیب الگوریتم بهینه‌سازی ذرات و الگوریتم ژنتیک. هشتمین کنگره ملی تازه‌های مهندسی برق و کامپیوتر ایران، تهران.

\bibitem{ref26} الوانی، سید مهدی و دانایی فرد، حسن (۱۳۸۴). تئوری نظم در بی‌نظمی و مدیریت. تهران: انتشارات صفار.

\end{thebibliography}
\end{RTL}
\
